\chapter{From Score to Velocity: Riemannian Flow Matching}
\label{ch:theory}

This chapter develops the generative machinery that replaces diffusion in
SigmaFlow. The treatment is selective: results that are standard and not
load-bearing for the argument are stated and cited, while two developments
--- the conditional objective and the geodesic construction on $\SOthree$ ---
are derived, because Chapters~\ref{ch:asymmetry} and \ref{ch:results} depend
on their precise form.

\section{Continuous normalising flows and the intractable marginal field}

Let $\Mfd$ be the state space and let $\vfield_t$ be a time-dependent vector
field on it. The associated flow $\psi_t$ solves
\begin{equation}
  \frac{\dif}{\dif t}\psi_t(x) \;=\; \vfield_t\bigl(\psi_t(x)\bigr),
  \qquad \psi_0(x) = x .
  \label{eq:flow-ode}
\end{equation}
Pushing a source distribution $p_0$ through $\psi_t$ produces a path of
distributions $p_t = (\psi_t)_\# p_0$, and this is a continuous normalising
flow \citep{chen2018neural}. Generative modelling reduces to finding a field
whose flow transports a tractable $p_0$ to the data distribution $p_1$.

If such a field is known, sampling is straightforward: draw $x_0 \sim p_0$ and
integrate \eqref{eq:flow-ode}. The difficulty is training. A field
$\vfield_t$ generates $p_t$ precisely when the pair satisfies the continuity
equation, $\partial_t p_t + \nabla\!\cdot\!(p_t\,\vfield_t) = 0$, and the
natural objective
\begin{equation}
  \Loss_{\mathrm{FM}}(\theta)
  = \E_{t,\,x \sim p_t}\bigl\|\,\pred(t,x) - \vfield_t(x)\,\bigr\|^2
  \label{eq:fm-loss}
\end{equation}
cannot be evaluated, because neither $p_t$ nor $\vfield_t$ is available in
closed form for any interesting $p_1$. Flow matching resolves this
\citep{lipman2023flow}.

\section{Conditional flow matching}
\label{sec:cfm}

The construction conditions on the data point. For each $x_1$, choose a
conditional path $p_t(\,\cdot \mid x_1)$ with $p_0(\cdot \mid x_1) = p_0$ and
$p_1(\cdot \mid x_1) = \delta_{x_1}$, generated by a conditional field
$\vfield_t(\cdot \mid x_1)$ that we may write down by hand. The marginal path
and its field are then
\begin{equation}
  p_t(x) = \int p_t(x \mid x_1)\, q(x_1)\, \dif x_1,
  \qquad
  \vfield_t(x) = \int \vfield_t(x \mid x_1)\,
  \frac{p_t(x \mid x_1)\, q(x_1)}{p_t(x)}\, \dif x_1 ,
  \label{eq:marginal-field}
\end{equation}
where $q$ is the data distribution. The second expression is exactly a
conditional expectation,
\begin{equation}
  \vfield_t(x) \;=\; \E\bigl[\,\vfield_t(x \mid X_1) \,\big|\, X_t = x\,\bigr] .
  \label{eq:marginal-as-cexp}
\end{equation}
This identity is the pivot of the whole method and is used again, for a quite
different purpose, in Chapter~\ref{ch:asymmetry}.

\begin{proposition}[Marginalisation]
\label{prop:marginalisation}
Define the conditional objective
\begin{equation}
  \Loss_{\mathrm{CFM}}(\theta)
  = \E_{t,\,x_1 \sim q,\; x \sim p_t(\cdot \mid x_1)}
    \bigl\|\,\pred(t,x) - \vfield_t(x \mid x_1)\,\bigr\|^2 .
  \label{eq:cfm-loss}
\end{equation}
Then $\nabla_\theta \Loss_{\mathrm{CFM}} = \nabla_\theta \Loss_{\mathrm{FM}}$,
so the two objectives have the same minimisers in $\theta$.
\end{proposition}

\begin{proof}[Proof sketch]
Expand both squared norms. The terms $\|\pred\|^2$ agree because both
expectations are over the same marginal $p_t$. The terms not involving
$\theta$ do not affect the gradient. The cross terms agree because, by
\eqref{eq:marginal-as-cexp} and the tower property,
$\E_{p_t}\langle \pred(t,x), \vfield_t(x)\rangle
 = \E_{p_t}\langle \pred(t,x), \E[\vfield_t(x\mid X_1)\mid X_t=x]\rangle
 = \E_{x_1, x}\langle \pred(t,x), \vfield_t(x \mid x_1)\rangle$.
\end{proof}

Two consequences follow, and they pull in different directions.

The practical one is that training becomes a regression onto a target written
down in closed form. Nothing is simulated: one samples $t$, samples a data
point, samples a point on its conditional path, evaluates
$\vfield_t(x \mid x_1)$ analytically, and takes a gradient step. Diffusion
models achieve something similar through denoising score matching
\citep{vincent2011connection}, but there the tractable target arises from a
Gaussian transition kernel, whereas here the path is chosen first and the
target follows from it.

The second consequence is subtler and is the one this thesis turns on. The
minimiser of \eqref{eq:cfm-loss} is not the conditional field --- that object
depends on $x_1$, which the network never sees --- but its conditional
expectation \eqref{eq:marginal-as-cexp} given the current state. This is
exactly the standard fact that an $L^2$ regression returns a conditional mean.
Its force depends entirely on how much $X_t$ reveals about $X_1$. If the state
nearly determines the data point, the conditional mean is close to a single
target and the regression is well posed. If it does not, the model is trained
towards an average of many incompatible targets, and that average may be a
poor member of the set it averages. Chapter~\ref{ch:asymmetry} shows that the
two factors of the docking state space sit at opposite ends of this spectrum.

In $\R^d$ the standard choice is the straight-line path
\begin{equation}
  x_t = (1-t)\,x_0 + t\,x_1,
  \qquad
  \vfield_t(x_t \mid x_0, x_1) = x_1 - x_0 ,
  \label{eq:linear-path}
\end{equation}
whose conditional trajectories are straight and whose target is constant along
each one. This is what SigmaFlow uses on the translational factor.

\section{The Riemannian generalisation}
\label{sec:riemannian}

On a Riemannian manifold $(\Mfd, g)$ several ingredients of the Euclidean
construction stop making sense, and it is worth being precise about which.

The interpolant $(1-t)x_0 + t x_1$ is not defined, because there is no
addition. The difference $x_1 - x_0$ is not defined either, and neither is the
statement that a velocity ``is'' a point of the space: a velocity at $x$ is an
element of the tangent space $T_x\Mfd$, and tangent spaces at different points
are different vector spaces that cannot be compared without extra structure.
Finally, the squared error $\|\pred - \vfield\|^2$ must be measured in the
metric at the point where both vectors live.

What survives is the part that matters. Proposition~\ref{prop:marginalisation}
rests only on bilinearity of an inner product and the tower property of
conditional expectation, both of which hold pointwise on a manifold, so the
conditional objective remains a valid surrogate \citep{chen2024riemannian}.
Three substitutions then carry the construction over: straight lines become
geodesics; differences become logarithmic maps; and the Euclidean norm in
\eqref{eq:cfm-loss} becomes the Riemannian metric at the relevant point,
\begin{equation}
  \Loss_{\mathrm{RCFM}}(\theta)
  = \E_{t,\,x_1,\,x_t}
    \bigl\|\,\pred(t,x_t) - \vfield_t(x_t \mid x_1)\,\bigr\|_{g(x_t)}^2 .
  \label{eq:rcfm-loss}
\end{equation}

The geodesic interpolant between $x_0$ and $x_1$ is
\begin{equation}
  x_t \;=\; \Exp_{x_0}\!\bigl(t \,\Log_{x_0}(x_1)\bigr),
  \label{eq:geodesic-interp}
\end{equation}
where $\Exp_x : T_x\Mfd \to \Mfd$ follows the geodesic leaving $x$ with a
given initial velocity and $\Log_x$ is its local inverse. Differentiating
\eqref{eq:geodesic-interp} and re-expressing the result at $x_t$ gives the
conditional field
\begin{equation}
  \vfield_t(x_t \mid x_1) \;=\; \frac{\Log_{x_t}(x_1)}{1-t} ,
  \label{eq:riem-target}
\end{equation}
which is the statement that at time $t$ one must travel the remaining geodesic
distance in the remaining time. Equation~\eqref{eq:riem-target} is the object
whose behaviour on $\SOthree$ drives Chapter~\ref{ch:asymmetry}.

\section{The rotation group}
\label{sec:so3}

$\SOthree = \{\Rot \in \R^{3\times 3} : \Rot^\top \Rot = I, \det \Rot = 1\}$
is a compact three-dimensional Lie group. Its Lie algebra $\sothree$ is the
space of skew-symmetric matrices, identified with $\R^3$ by the hat map
\begin{equation}
  \skewm{\omega} =
  \begin{pmatrix} 0 & -\omega_3 & \omega_2\\
                  \omega_3 & 0 & -\omega_1\\
                  -\omega_2 & \omega_1 & 0 \end{pmatrix},
  \qquad \omega \in \R^3 .
\end{equation}
With the bi-invariant metric $\langle A, B\rangle = \tfrac12\Tr(A^\top B)$ the
exponential map is Rodrigues' formula and its inverse is available in closed
form for $\theta \in (0,\pi)$:
\begin{equation}
  \Exp(\skewm{\omega}) = I + \frac{\sin\theta}{\theta}\skewm{\omega}
    + \frac{1-\cos\theta}{\theta^2}\skewm{\omega}^2,
  \qquad
  \theta = \|\omega\| ,
  \label{eq:rodrigues}
\end{equation}
\begin{equation}
  \Log(\Rot) = \frac{\theta}{2\sin\theta}\bigl(\Rot - \Rot^\top\bigr),
  \qquad
  \theta = \arccos\!\Bigl(\frac{\Tr \Rot - 1}{2}\Bigr) .
  \label{eq:so3-log}
\end{equation}
The choice of a \emph{bi-invariant} metric is not cosmetic. It makes left and
right translation isometries, so distances are unchanged by re-expressing both
arguments in another frame; it makes the geodesics through the identity
exactly the one-parameter subgroups $t \mapsto \Exp(t\skewm{\omega})$, which is
what puts \eqref{eq:rodrigues} in closed form; and it makes the Haar measure
the Riemannian volume, so ``uniform on $\SOthree$'' is unambiguous. Every
closed form in this chapter depends on it, and it is available because
$\SOthree$ is compact.

The geodesic distance between $\Rot_a$ and $\Rot_b$ is
$\|\Log(\Rot_a^\top \Rot_b)\|$, taking values in $[0,\pi]$. The logarithm is
branch-ambiguous at $\theta = \pi$, the cut locus: a rotation by exactly $\pi$
about an axis $n$ is equally a rotation by $-\pi$ about the same axis, and
$\Log$ must choose. This is not a rare edge case here. Under the Haar measure
the angle has density proportional to $1 - \cos\theta$, which is
\emph{increasing} on $[0,\pi]$, so uniformly sampled rotations concentrate
towards large angles: roughly one draw in ten exceeds $\Ang{170}$. Since the
rotational source of \S\ref{sec:sigmaflow-objective} is exactly Haar, a tenth
of all training pairs begin near the cut locus, and the numerical treatment of
\eqref{eq:so3-log} there is a correctness requirement rather than a
refinement. \S\ref{sec:log-defect} reports a defect of exactly this kind.

\subsection{Trivialisation, fixed explicitly}
\label{sec:trivialisation}

Because $\SOthree$ is a group, a tangent vector at $\Rot$ may be transported
to the identity in two inequivalent ways, and the resulting conventions are
easy to conflate. Writing a curve's velocity as
\begin{equation}
  \dot{\Rot} = \Rot\,\skewm{\Omega}
  \;\Longleftrightarrow\;
  \skewm{\Omega} = \Rot^\top \dot{\Rot}
  \qquad\text{(body frame; \emph{left}-trivialised)}
  \label{eq:body-frame}
\end{equation}
differs from
\begin{equation}
  \dot{\Rot} = \skewm{w}\,\Rot
  \;\Longleftrightarrow\;
  \skewm{w} = \dot{\Rot}\,\Rot^\top
  \qquad\text{(spatial frame; \emph{right}-trivialised)},
  \label{eq:spatial-frame}
\end{equation}
and the two are related by the adjoint action,
\begin{equation}
  \skewm{\Omega} \;=\; \Rot^\top \skewm{w}\, \Rot
  \;=\; \Ad_{\Rot^{-1}} \skewm{w} .
  \label{eq:adjoint}
\end{equation}
This thesis uses the \emph{body-frame} convention \eqref{eq:body-frame}
throughout, following \citet{marsden1999introduction}. Every rotational
velocity, target and network output is a body-frame quantity, and the
integrator correspondingly right-multiplies. \S\ref{sec:frame-defect}
documents a defect that arose from mixing \eqref{eq:body-frame} and
\eqref{eq:spatial-frame}, and the terminological trap that produced it.

\section{The SigmaFlow probability path and objective}
\label{sec:sigmaflow-objective}

Specialising \eqref{eq:geodesic-interp} and \eqref{eq:riem-target} to the two
factors of $\Mfd$ gives the construction implemented in SigmaFlow. On each
translational factor, with $x_0 \sim \mathcal{N}(0, I_3)$ in pocket-centred
coordinates,
\begin{equation}
  x_t = (1-t)x_0 + t\,x_1,
  \qquad
  \vfield_t^{\mathrm{trans}} = x_1 - x_0
  \;=\; \frac{x_1 - x_t}{1-t} .
  \label{eq:sf-trans}
\end{equation}
On each rotational factor, with $\Rot_0 \sim \Haar$ drawn from the Haar
measure,
\begin{equation}
  \Rot_t \;=\; \Rot_0 \Exp\!\bigl(t\,\Log(\Rot_0^\top \Rot_1)\bigr),
  \qquad
  \vfield_t^{\mathrm{rot}} \;=\; \frac{\Log(\Rot_t^\top \Rot_1)}{1-t} .
  \label{eq:sf-rot}
\end{equation}
Both expressions in \eqref{eq:sf-trans} are equal, and the corresponding
identity for \eqref{eq:sf-rot} is proved in Chapter~\ref{ch:asymmetry}; the
implementation uses whichever form is better conditioned. The right-hand
forms, involving $x_t$ and $\Rot_t$, express the target as ``the remaining
distance divided by the remaining time'' and are singular at $t=1$; the
left-hand forms, involving only the endpoints, are not.

The product structure of $\Mfd$ makes the extension to $N$ fragments
immediate. On a Riemannian product the metric is the direct sum of the factor
metrics, geodesics are products of factor geodesics, and $\Exp$ and $\Log$ act
componentwise. A path constructed independently on each factor is therefore a
legitimate path on $\Mfd$, and the conditional field is the concatenation of
the factor fields. Since the factors are treated as independent, the objective
is a weighted sum of the two factor-wise losses,
\begin{equation}
  \Loss(\theta) \;=\;
  \E_{t \sim \mathcal{U}[0,1]}\Bigl[
     \lambda_{\mathrm{trans}}\bigl\|\pred^{\mathrm{trans}} - \vfield_t^{\mathrm{trans}}\bigr\|^2
   + \lambda_{\mathrm{rot}}\bigl\|\pred^{\mathrm{rot}} - \vfield_t^{\mathrm{rot}}\bigr\|^2
  \Bigr],
  \label{eq:sf-objective}
\end{equation}
with the weights inherited unchanged from SigmaDock's production
configuration. Sampling integrates the learned field by explicit Euler,
using the exponential map as the retraction on the rotational factor:
\begin{equation}
  x^{(k+1)} = x^{(k)} + \pred^{\mathrm{trans}}\,\Delta t,
  \qquad
  \Rot^{(k+1)} = \Rot^{(k)}\Exp\!\bigl(\pred^{\mathrm{rot}}\,\Delta t\bigr) .
  \label{eq:euler}
\end{equation}
The right-multiplication in \eqref{eq:euler} is the concrete expression of the
body-frame convention fixed in \S\ref{sec:trivialisation}.

\section{Score versus velocity}
\label{sec:score-vs-velocity}

The substitution changes what the network means, not what it outputs. Both
models emit the per-atom field \eqref{eq:force-field} and reduce it by
\eqref{eq:reduce-trans}--\eqref{eq:reduce-rot}. SigmaDock interprets the
result as a score $\nabla \log p_\tau$, whose magnitude must be rescaled by
noise-level-dependent factors before it can be used in an update; SigmaFlow
interprets it as a velocity, which is already in the units of the state and
requires no rescaling. The entire scalar apparatus that converts scores into
displacements is therefore deleted rather than modified
(\S\ref{sec:minimal-change}).

The two formulations are relatives. A diffusion model admits a deterministic
probability-flow ODE whose velocity is an affine function of the score
\citep{song2021sde}, so diffusion may be viewed as flow matching along a
particular, Gaussian, probability path \citep{lipman2023flow}. What
distinguishes them in practice is that flow matching makes the path an
explicit design choice, and that sampling is deterministic.

\section{The source distribution as a modelling choice}
\label{sec:source-theory}

One consequence of that freedom deserves emphasis, because
Chapter~\ref{ch:asymmetry} turns on it and \S\ref{sec:source-results} exploits
it. Under diffusion the source is not chosen: it is whatever the forward
noising process converges to, an isotropic Gaussian on $\R^3$ and the
uniform measure on $\SOthree$. Under flow matching the only requirements on
$p_0$ are that it be tractable to sample and that the conditional path connect
it to $p_1$. Any such $p_0$ is admissible, including one that depends on the
conditioning context $\ctx$.

This is a genuinely flow-specific degree of freedom, and it is not merely a
convenience. The quality of the learning problem depends on how much the
source already knows. SigmaFlow's translational source is centred on the
binding pocket and is therefore informative; its rotational source is Haar and
carries no orientational information whatsoever. Chapter~\ref{ch:asymmetry}
argues that this difference has consequences that the geometry makes precise.
